\section{ROSE Installation}

See http://rosecompiler.org/ROSE\_HTML\_Reference/installation.html.

\subsection{Software/Hardware Requirements and Options}
\label{Requirements_Installation_Testing}

See http://rosecompiler.org/ROSE\_HTML\_Reference/installation.html

%%    ROSE is general software and we ultimately hope to remove any specific software
%% and hardware requirements.  However, our goal is to be specific about where and 
%% how ROSE is developed and where it is regularly tested.
%% % ROSE has been developed initially on the Sun workstations and later on Linux.
%% % Since ROSE is written in C++, it requires a C++ compiler. Some optional features
%% % within ROSE have some additional software requirements.

\subsubsection{Required Hardware/Operating System}

See http://rosecompiler.org/ROSE\_HTML\_Reference/installation\_hardware.html

%% % Irrelevant details.
%% %The ROSE C/C++ frontend, legacy frontend (software from the Edison Design Group),
%% %is released as source code or binary; all other language support in
%% %ROSE is fully released as source code. legacy frontend has addressed portability
%% %issues for their C++ frontend and it is available on nearly all
%% %platforms (see {\tt www.legacy frontend.com} for details); so legacy frontend is not a barrier
%% %to portability for us.
%% 
%% If you have a specific requirement for ROSE to be ported to a target platform
%% please let us know.
%% % At a later point we will address portability of ROSE onto other platforms:

\subsubsection{Software Requirements}

See http://rosecompiler.org/ROSE\_HTML\_Reference/installation\_prerequisites.html

%% %This is too verbose. I merged the text into the itemized list to quickly
%% %get to the point.
%% %   You will require a C++ compiler to compile ROSE, and some of the test
%% %cases also require a C compiler.
%% % We use the GNU g++ compiler most often, but the Intel compiler (version 9.1) will also work.
%% %Present development work is done on Intel/Linux platforms
%% %using the GNU g++ 3.3.x, 3.4.x, and 4.x; and the Intel compilers.
%% %using the GNU g++ 4.0.x to 4.4.x; and the Intel compilers.
%% %GNU compilers older
%% %than 3.3.x are not supported within the ROSE project, but in some cases they might work 
%% %(g++ 3.2.x is likely to work, but g++ 2.96 will not work).
%% 
%% %{\bf Required Software:} \\
%% %  The following software is required in order to build and use ROSE 
%% \begin{itemize}
%%    \item {\bf ROSE}: \\
%%      There are three distributions of ROSE: the {\it Distribution Version} for users
%%     (typical), the {\it External Development Version} for advanced users and collaborators, 
%%      and the {\it Internal Development Version}
%%    % The exact requirements are listed in the {\tt ROSE/ChangeLog} (including version numbers for each release of ROSE).
%%      \begin{itemize}
%%        \item {\bf Distribution Version} \\
%%        Provided as a tared and compressed file named ROSE-\VersionNumber.tar.gz.
%%        It can be obtained from \htmladdnormallink{https://outreach.scidac.gov/projects/rose}{https://outreach.scidac.gov/projects/rose/}. 
%%        This is the most typical way that users will see and work with ROSE. But it is less
%%        up-to-date compared to development versions.
%% 
%%        \item {\bf External Development Version:git} \\
%%        You can obtain it by typing \textit{git clone
%%        http://www.rosecompiler.org/rose.git} 
%% 
%%        %It is available from
%%        %\htmladdnormallink{http://www.rosecompiler.org/rose.git}{http://www.rosecompiler.org/rose.git}.
%%        %This is a clone from our internal ROSE repository. 
%%        It enables people to have the quickest access to the most recent new features in ROSE.  
%%        The external git repository is synchronized with the
%%        internal repository once its master branch is updated. 
%%        It is highly recommended to be used by collaborators who regularly contribute back to us. 
%%        %Several branches also exist to accept contributions from external collaborators.
%% 
%% 
%%        \item {\bf External Development Version: subversion} \\
%%        It is available from the SciDAC Outreach Center's subversion repository:
%%        \htmladdnormallink{https://outreach.scidac.gov/projects/rose}{https://outreach.scidac.gov/projects/rose/}.
%%        It contains only a subset of ROSE (excluding the legacy frontend part essentially) of the internal developer version. 
%%        We are trying to gradually phase out the subversion repository to
%%        facilitate merging external contributions.
%% 
%%        \item {\bf Internal Development Version: git repositories} \\
%%        Internally, we maintain two git repositories internally, one for ROSE and the
%%        other for legacy frontend source code. The legacy frontend repository is linked to the ROSE
%%        repository as a git submodule. 
%%        The details of getting and building this version are located in the
%%        Developer Guide. %\ref{gettingStarted:DeveloperInstructions}.
%%        This distribution is intended only for ROSE development team
%%        and external developers with access to our internal network file
%%        system.  
%%       % The development versions are what are found in the ROSE software
%%       % repositories and have additional software requirements: subversion, JDK, MrDocs, LaTeX, {\it etc}.
%%      \end{itemize}
%%     % ROSE is delivered as either a development version (svn) or as a distribution for
%%     % users. The (typical) user version comes either with legacy frontend source or binary distribution.
%% 
%%   \item {\bf git}: Required if you work with development versions of ROSE
%%   using git repositories.
%% 
%%   \item {\bf subversion}: Required if you work with development versions of ROSE
%%   using subversion repositories.
%% 
%%   \item {\bf wget}: Required for all external distribution/development
%%   versions. \\
%%   Only the internal git development version has access to legacy frontend source code.
%%   For all other cases,  wget is used to automatically download prebuilt legacy frontend binaries matching
%%   your platform, if the configure script detects that your platform is supported. 
%% 
%%    \item {\bf g++} : version 4.0.x to 4.4.x (Support for GCC 4.5 and higher
%%    versions is still ongoing work).
%% 
%%    \item {\bf gfortran}: version 4.2.x to 4.4.x. Required if you want to
%%    handle Fortran input.
%% 
%%      versions) right now due to resource limits. 
%% 
%%     \item {\bf Autoconf} : version $>=$ 2.59. Needed \emph{ONLY} for development versions. \\
%%      Autoconf is an extensible package of M4 macros that produce shell scripts to automatically configure software source code packages. 
%% 
%%    \item {\bf Automake} : version $>=$ 1.96. Needed \emph{ONLY} for development versions. \\
%%      The build system generates platform-specific build files compliant with
%%      the GNU Coding Standards.
%% 
%%    \item {\bf GNU Libtool}: version $>=$1.5.6.  Needed \emph{ONLY} for development versions. 
%% 
%%    \item {\bf GNU Flex}: Needed \emph{ONLY} for development versions. \\
%%    Flex is used to generate annotation\/OpenMP directive scanners in ROSE. 
%% 
%%    \item {\bf GNU Bison}: Needed \emph{ONLY} for development versions. \\
%%    Yacc is used to generate some annotation\/OpenMP directive parsers in ROSE. 
%%
%%    \item {\bf MrDocs} : Needed \emph{ONLY} for development versions. \\
%%           The API reference is generated using MrDocs, so ROSE developers that want
%%           to regenerate it need the MrDocs binary available in their PATH. This is
%%           not required for ROSE users, since all documentation is included in the
%%           ROSE distribution.
%%  
%%    \item {\bf ghostscript}: Needed \emph{ONLY} for development versions. \\
%%    ps2pdf is need  to generate ROSE documentations.
%% 
%%    \item {\bf DOT (GraphViz)} : Required since many ROSE tests generate
%%    AST/CFG graph in dot format.\\
%%           ROSE uses DOT for generating graphs of ASTs, Control Flow, etc.
%%           Visit {\it www.graphviz.org} for details and to download software.
%%           An example showing the use of the DOT to build graphs is in the ROSE Tutorial.
%%           There are no ROSE-specific configure options to use dot; it must only be 
%%           available within your path.
%%  
%%    \item {\bf LaTex: texlive}: Needed \emph{ONLY} for development versions. \\
%%      LaTeX is used for a significant portion of the ROSE documentation. LaTex is 
%%      included on most Unix systems. There are no ROSE specific configure options 
%%      to use LaTeX; it must only be available within your path.
%% 
%% \end{itemize}
%% 
%% {\bf Optional Software:} \\
%% %ROSE does not presently require any special software, but more
%% More functionality within ROSE is available if one has additional (freely available) software installed:
%% \begin{itemize}
%%    \item {\bf libxml2-devel }: Optional for some features. \\
%%          Several optional features of ROSE need to handle XML files.
%% 
%% 
%%  \item {\bf SQLite} : \\
%%           ROSE users can store persistent data across separate compilation of files by
%%           storing information in an SQLite database.  This is used by several features
%%           in ROSE (call-graph generation, for example) and may be used directly by the
%%           user for storage of user-defined analysis data.  Such database support is
%%           one way to handle global analysis (the other way is to build the whole
%%           application AST). Visit {\it www.sqlite.org} for details and to download
%%           software. %Details are in section \ref{gettingStarted:Database}.
%%           An example showing the use of the ROSE database mechanism is in the ROSE
%%           Tutorial. Use of SQLite requires special ROSE configuration options (so that
%%           the SQLite library can be added to the link line at compile time).  See ROSE
%%           configuration options for more details ({\tt configure --help}).
%%    \item {\bf mpicc} : \\
%%           mpicc is a compiler for MPI development. If ROSE is configures with MPI
%%           enabled, one can utilize features in ROSE that allow for distributed
%%           parallel AST traversals.
%%    \item {\bf Haskell (ghc) Compiler}: \\
%%           ROSE supports Haskell bindings so that ROSE tools can be written in Haskell.
%%           ROSE does not parse of permit source-to-source analysis and transformation of
%%           Haskell langauge.  So we support Haskell but only in this specific way.
%%           See directions below for how to install the Haskell (ghc) compiler.
%% 
%% \end{itemize}
%% 
%% %{\bf Required Software for ROSE Demos:} \\
%% %     The ROSE demos make use of additional external software which we also find useful for
%% %     internal development.
%% %\begin{itemize}
%% %\end{itemize}
%% 
%% 
%% 
%% % \subsection{Installation and Testing }
%% 
%% %\subsection{How to Build ROSE}
%% 
%% \commentout{
%% %   This section addresses the building and installation of ROSE.  
%% There are two versions of ROSE supported: the 
%% {\it Distribution Version} for users (typical) and the {\it Development Version}
%% (intended only for ROSE development team), which is what is found in the ROSE software
%% repository and has additional software requirements (MrDocs, LaTeX,
%% etc.).  The exact requirements are listed in the {\tt ROSE/ChangeLog} (including version 
%% numbers for each release of ROSE).
%% \begin{itemize}
%%    \item {\bf Distribution Version} \\
%%          Provided as a tared and compressed file in the form,
%%          ROSE-\VersionNumber.tar.gz.  This is the most typical way that users will
%%          see and work with ROSE.  Instructions are also located in the ROSE/README file.
%% 
%%    \item {\bf Development Version} \\
%%          Only available directly from the Subversion (SVN) repository. The details of building this
%%          version are located in the Appendix \ref{gettingStarted:DeveloperInstructions}.
%% \end{itemize}  
%% To simplify the descriptions of the build process, we define:
%% \begin{itemize} 
%%    \item {\bf Source Tree} \\
%%          Location of source code directory (there is only one source tree).
%%    \item {\bf Compile Tree} \\
%%          Location of compiled object code and executables. There can be many compile trees
%%          representing either different: configure options, compilers used to build ROSE
%%          and ROSE translators, compilers specified as backends for ROSE (to compile ROSE
%%          generated code), or architectures.
%%    \item {\bf Install Tree} \\
%%          Location of compiled object code when copied to a clean installation directory.
%%          This is commonly used in order to refer to a project from outside, i.e. to 
%%          use the provided interfaces. 
%% \end{itemize}
%% 
%% We {\em strongly} recommend that the {\bf Source Tree} and the {\bf Compile Tree} be
%% different. This avoids many potential problems with the {\em make clean} rules.
%% Note that the {\bf Compile Tree} will be the same as the {\bf Source Tree}
%% if the user has {\em not} explicitly generated a separate directory in which to
%% run {\em configure} and compile ROSE. If the {\bf Source Tree} and {\bf Compile Tree} 
%% are the same, then there is only one combined {\bf Source/Compile Tree}.
%% Alternatively, numerous different {\bf Compile Trees} can be used from 
%% a single {\bf Source Tree}.  More than one {\bf Compile Tree} allows
%% ROSE to be generated on different platforms from a single source 
%% (either a generated distribution or checked out from SVN).  ROSE is developed 
%% and tested internally using separate {\bf Compile Trees}.
%% }
%% 
%% \commentout{
%% \subsection{Summary of Build Process}
%% \label{gettingStarted:SummaryInstructions}
%%    More specific information is available in the subsections below 
%% (subsection \ref{gettingStarted:UserInstructions} and 
%% subsection \ref{gettingStarted:DeveloperInstructions}) and various README files in the top of the source tree, but in summary the steps are:
%% % The general steps are (but use the more details instructions below, depending 
%% % on the version of ROSE that you have):
%% \begin{itemize}
%%    \item Type {\tt \<Source Tree\>/configure;} \\
%%          This can be done in a separate directory or in the source directory.
%%          We suggest using a separate directory, but either will work. The 
%%          {\bf Source Tree} may be specified with either an absolute or relative 
%%          path.
%%    \item type {\tt make} \\ to build the source code (you may also use the parallel make
%%          option (e.g. {\tt make -j4} to run make using 4 processes).
%%    \item Type {\tt make check} \\ to run internal tests to test your build.
%%    \item Type {\tt make install} \\ to install the build library for more convenient use.
%%          (default is {\tt /usr/local}, use {\tt --prefix=`pwd'} to specify the current directory).
%%    \item Type {\tt make documentation} \\ to build both html and LaTeX documentation.
%% \end{itemize}
%% }
%% 
\subsection{Using Insure++}
\label{gettingStarted:insure}

   Insure++ integration is experimental and not currently wired into the CMake build.
Use external tooling if you need this kind of analysis support.

% I don't know if this actually works since I use the setup script.
% If you need to specify the location of {\tt insure++} then use {\em --with-insure=<path>}.


\subsection{Building ROSE From a Distribution (ROSE-\VersionNumber.tar.gz)}
\label{gettingStarted:UserInstructions}

ROSE is no longer distributed as a tarball. Use the Git repositories
instead.

See http://rosecompiler.org/ROSE\_HTML\_Reference/installation.html

%% To simplify the descriptions of the build process, we define:
%% \begin{itemize} 
%%    \item {\bf Source Tree} \\
%%          Location of source code directory (there is only one source tree).
%%    \item {\bf Compile Tree} \\
%%          Location of compiled object code and executables. There can be many compile trees
%%          representing either different: configure options, compilers used to build ROSE
%%          and ROSE translators, compilers specified as backends for ROSE (to compile ROSE
%%          generated code), or architectures.
%% \end{itemize}
%% 
%% We {\em strongly} recommend that the {\bf Source Tree} and the {\bf Compile Tree} be
%% different. This avoids many potential problems with the {\em make clean} rules.
%% Note that the {\bf Compile Tree} will be the same as the {\bf Source Tree}
%% if the user has {\em not} explicitly generated a separate directory in which to
%% run {\em configure} and compile ROSE. If the {\bf Source Tree} and {\bf Compile Tree} 
%% are the same, then there is only one combined {\bf Source/Compile Tree}.
%% Alternatively, numerous different {\bf Compile Trees} can be used from 
%% a single {\bf Source Tree}.  More than one {\bf Compile Tree} allows
%% ROSE to be generated on different platforms from a single source 
%% (either a generated distribution or checked out from svn/git repos).  ROSE is developed 
%% and tested internally using separate {\bf Compile Trees}. \\
%% 
%% 
%%    The process for building ROSE from a ROSE {\em Distribution Version} is the same as for
%%    most standard software distributions. Specific details for particular
%%    operating system distributions can be found in various README files in the top of the source tree. In
%%    general:
%% \begin{enumerate}
%%      \item Untar ROSE. \\
%%            Type {\tt tar -zxf ROSE-\VersionNumber.tar.gz} to untar the ROSE distribution.
%%      \item Build a separate compile tree. \\
%%            Type {\tt mkdir compileTree} to build a location for the object files and
%%            documentation (use any name you like for this directory).
%%      \item Change directory to the new compile tree directory. \\
%%            Type {\tt cd compileTree; }. This changes the current directory to the newly
%%            created directory.
%%      \item Add JVM environment variables for OFP if your platform is Linux. \\
%%            For example:\\
%%            {\tt export JAVA\_HOME=/usr/apps/java/jdk1.5.0\_11} \\
%%            {\tt export LD\_LIBRARY\_PATH=\$JAVA\_HOME/jre/lib/i386/server:\$LD\_LIBRARY\_PATH}
%% 
%%      the path in front of any other paths to avoid some old versions of
%%            For example: For Linux, 
%% 
%% 
%%      \item Run the {\tt configure} script. \\
%%            Type {\tt \{AbsoluteOrRelativePath\}/configure
%%            to run the ROSE {\tt configure} script.  The path to the configure script 
%%            may be either relative or absolute. The prefix option on the configure 
%%            command line is only required if you run {\tt make install} (suggested), because 
%%            the default location for installation is {\tt \//usr\//local} and most users don't
%%            have permission to write to that directory. This is common to all projects that
%%            use the standard GNU Makefile layout for installed artifacts. As of ROSE-0.8.9a, the
%%            default setting for the install directory (prefix) is the build tree.
%%            For more on ROSE configure options, see section \ref{gettingStarted:configureOptions}.
%%            {\em Note: If configure is not present, you may have a development version. 
%%            In that case, see the next section
%%            (\ref{gettingStarted:UserInstructionsForDevelopmentVersion}) for installation instructions.}
%% 
%%      \item Run {\tt make}. \\
%%            Type {\tt make} to build all the source files. See details of running 
%%            {\tt make} in parallel in section \ref{gettingStarted:parallelMake}.
%% 
%%      \item To test ROSE (optional). \\
%%            Type {\tt make check} to test the ROSE library against a collection of test codes.
%%            See details of running make in parallel \ref{gettingStarted:parallelMake}.
%% 
%%      \item To install ROSE, type {\tt make install}. \\
%%            Installation is optional, but suggested. Users can simplify their use of ROSE 
%%            by using it from an installed version of ROSE.  This permits compilation 
%%            using a single include directory and the specification of only two libraries.
%%            See details of installing ROSE in section \ref{gettingStarted:installation}.
%% 
%%      \item Testing the installation of ROSE (optional). \\
%%            To test the installation and the location where ROSE is installed, against a
%%            collection of test codes (the application examples in {\tt ROSE/tutorial}), 
%%            type {\tt make installcheck}.
%%            A sample {\tt makefile} is generated.%; see section \ref{gettingStarted:compilingTranslator}.
%% \end{enumerate}

\subsection{Building ROSE from a Development Version (from SVN or GIT)}
\label{gettingStarted:UserInstructionsForDevelopmentVersion}

See http://rosecompiler.org/ROSE\_HTML\_Reference/installation.html

%% Building ROSE from an internal or external development version is very
%% similar to building it from a distribution. The major difference is that
%% for development versions, there is no configure script available . 
%% 
%% You have to type {\tt ./build} in the source tree to 
%% generate the build configuration files. Once this is done, the rest
%% steps are the same as those of building a distribution version.

\subsection{TroubleShooting the ROSE Installation}

% DQ (11/12/2008): This is where we need to accumulate the symptoms of what goes wrong.

   There are a number of famous ways to screw up your installation of ROSE.
\begin{enumerate}
   \item If CMake cannot locate the JVM needed for the Open Fortran Parser (OFP),
   ensure {\tt JAVA\_HOME} points to a JDK and {\tt LD\_LIBRARY\_PATH} includes the
   JVM server library (e.g., {\em libjvm.so}).

   \item If you see stale generated-file or build errors after a large update,
   create a fresh build directory and reconfigure with CMake.
\end{enumerate}



\subsection{ROSE CMake Options}
\label{gettingStarted:configureOptions}

     For CMake configuration options and documentation, run:
\begin{verbatim}
cmake -S . -B build -LH
\end{verbatim}
     or use {\tt ccmake} / {\tt cmake-gui} for an interactive view.

%%      A few example configure options are:
%% \begin{itemize}
%%    \item Minimal configuration \\
%%          This will configure ROSE to be compiled in the current directory (separate from
%%          the {\bf Source Tree}).  The installation (from {\tt make install}) will be placed in
%%          {\tt /usr/local}.  Most users don't have permission to write to this directory,
%%          so we suggest always including the {\em prefix option} (e.g. {\tt
%%          --prefix=`pwd`}).
%%    \item Minimal configuration (preferred) \\
%%          Configure in the current directory so that installation will also happen in the
%%          current  directory (a {\tt install} subdirectory will be built).
%%    \item Debugging \\
%%          Note that all possible debug modes are turned on by default. These include:
%%          \begin{itemize}
%%             \item {\tt -g } (compiler debug flag), and
%%             \item {\tt -Wall} (compiler warnings flag).
%%          \end{itemize}
%%          Additional modes are possible, one of which is now tested within ROSE internal
%%          release level testing:
%%          \begin{itemize}
%%             \item {\tt -D\_GLIBCXX\_DEBUG} (stl and other stdlib debug flag (used)),
%%             \item {\tt -D\_GLIBCXX\_DEBUG\_PEDANTIC"} (another stl debug flag we consider excessive (not used)).
%%          \end{itemize}
%%          change the size of iterators). 
%%          The use of these tests are eqvivalent to the configure line: \\
%%          Configure as above, if you like, but with debugging and warnings turned on.
%%          Note that if you compile ROSE with \verb!_GLIBCXX_DEBUG!, you should also compile
%%          the default is to turn on the use of {\tt -g } and {\tt -Wall} only.
%% %   \item Turning on compiler debugging options (\textbf{strongly recommended}) \\
%% %         Configure as above, but with debugging and warnings turned on.
%% %          Note that if you compile ROSE with \verb!_GLIBCXX_DEBUG!, you should also compile
%%    \item Adding Fortran support \\
%%          The Open Fortran Parser will also be enabled, allowing ROSE to process Fortran
%%          code. The JVM runtime must be available so the OFP frontend can run.
%%    \item Adding SQLite support \\
%%          Configure as above, but permit use of SQLite database for storage of analysis
%%          results between compilation of separate files (one type of support in ROSE for
%%          global analysis).
%%    \item Adding parallel distributed memory analysis support (using MPI) \\
%%          Configure as above, but with MPI and OpenMP support for ROSE to run AST traversals
%%          in parallel (distributed and shared memory).
%%    \item Adding IDA Pro support \\
%%          The binarysql flag allows ROSE to read a binary file previously stored as a sql file (e.g. fetched from
%%          IDA Pro). 
%%    \item Additional Examples \\ 
%%          More detailed documentation on configure options can be found by typing 
%%          {\tt configure --help}, or see figure
%%          \ref{gettingStarted:configureHelpOptionOutputPart1} for complete
%%          listing.
%% \end{itemize}



%\fixme{It is suggested by TID that we move the figures to this location.
%       Checkout if this is possible.}

\subsection{Running Parallel Builds}
\label{gettingStarted:parallelMake}

Use CMake's parallel build support, e.g.:
\begin{verbatim}
cmake --build build --parallel 16
\end{verbatim}

%%      ROSE uses general {\tt Makefiles} and is not dependent on {\em GNU Make}.
%%      However, {\em GNU Make} has an option to permit compilation in parallel and 
%%      we support this. Thus you may use {\tt make} with the {\tt -j<n>} option if 
%%      you want to run {\tt make} in parallel (a good value for {\tt n} is typically 
%%      twice the number of processors in your computer).  We have paid special 
%%      attention to the design of the ROSE {\em makefiles} to permit parallel 
%%      {\tt make} to work; we also use it regularly within development work.

\subsection{Installing ROSE}
\label{gettingStarted:installation}

See http://rosecompiler.org/ROSE\_HTML\_Reference/installation.html

%%      Installation (using {\tt make install}) is optional, but suggested. Users can
%% simplify their use of ROSE by using it from an installed version of ROSE.  This permits
%% compilation using a single include directory and the specification of only two libraries, 
%% as in:
%% \begin{verbatim}
%%      g++ -I{\<install dir\>/include} -o executable executable.C 
%%          -L{\<install dir\>/lib} -lrose -lm $(RT_LIBS)
%% \end{verbatim}
%% % $ reset emacs highlighting
%% 
%% See the example makefile in \\
%% {\tt ROSE/exampleTranslators/documentedExamples/simpleTranslatorExamples/exampleMakefile} \\
%% in Section \ref{gettingStarted:compilingTranslator}
%% for exact details of building a translator on your machine (setup by configure and tested
%% by {\tt make installcheck}).  Note that the tutorial example codes are also tested
%% by {\tt make installcheck} and the {\tt example\_makefile} there can also serve as
%% an example.
%%
%%      The default install prefix is {\tt /usr/local}. Only {\em root} has write privileges to that directory, so you
%% will likely get an error if you have not overridden the default value with a new
%% location.  To change the location, you need to have used the 
%% {\tt --prefix=\{install\_dir\}} to run the {\tt configure} script.  You can 
%% rerun the {\tt configure} script without rebuilding ROSE.
%% 
%% %\fixme{ROSE-0.8.9a now sets the default install path (prefix) to the build tree. The
%% %    documentation does not yet really reflect this detail.  Also since the build process
%% %    writes to the install directory we have to explain this in the documentation.}
%% 
%% % This MPI support section causes more confusion than benefits. I hide it
%% % from most users  Liao 4/14/2011
%% %\subsection{MPI Support}
%% %     ROSE supports the use of MPI for parallel distributed memory program analysis, a
%% %research focus within the ROSE project.  To support this use the 
%% %{\tt --with-parallel_ast_traversal_mpi}
%% %option on the configure command line.  If you get the following message:
%% %%{\footnotesize
%% %\begin{verbatim}
%% %configuration file /home/<user name>/.mpd.conf not found
%% %A file named .mpd.conf file must be present in the user's home
%% %directory (/etc/mpd.conf if root) with read and write access
%% %only for the user, and must contain at least a line with:
%% %MPD_SECRETWORD=<secretword>
%% %One way to safely create this file is to do the following:
%% %  cd $HOME
%% %  touch .mpd.conf
%% %  chmod 600 .mpd.conf
%% %and then use an editor to insert a line like
%% %  MPD_SECRETWORD=mr45-j9z
%% %into the file.  (Of course use some other secret word than mr45-j9z.)
%% %\end{verbatim}
%% %%}
%% %% $ reset emacs highlighting
%% %Then follow the directions to build the {\tt .mpd.conf} file.  The use of the 
%% %MPI configure option will allow additional code in ROSE to be compiled and 
%% %additional tests to be run.
%% 
%% \subsection{Haskell Support}
%%    Haskell (ghc) is available at {\it http://www.haskell.org/ghc/}.
%% To install ghc (from the binary distribution) run "configure" and 
%% "make install". Specifically (for example): \\
%% {\tt configure --prefix=/home/dquinlan/local/Haskell/ghc-6.10.4\_install}
%% and
%% {\tt make install}.
%% (This assumes that the path in your .bashrc includes ghc, e.g.: \\
%%    {\tt /home/dquinlan/local/Haskell/ghc-6.10.4\_install/bin}
%% 
%% That should be all that is required.  ROSE will detect
%% the existence of "ghc" at "configure" time and setup
%% what is required internally.  To turn off the use of Haskell support, 
%% use the configure option {\em --with-haskell=no}.
%% 
%% Note that we enforce and require version {\em 6.10.x} of ghc, this is
%% because we assume some default modules that have been removed in more 
%% recently released versions of ghc.
%% 
%% \subsection{Testing ROSE}
%%      A set of test programs is available.  %More details are in section \ref{testing}.
%% % written about this in later versions of the manual.  
%% Type {\tt make check} to run your build version of ROSE using these test codes.  
%% Several years of contributed bug reports and internal test codes have been accumulated 
%% in the {\tt ROSE/tests} directory.
%% 
%% Extra tests are available for development versions of ROSE. ROSE developers
%% are highly recommended to run {\tt make dist} and {\tt make distcheck} to make
%% sure that the modified development versions can be used to create functioning
%% distributions.
%% 
\subsection{Getting Help}
%     You may use the following mailing list to ask for help from the ROSE development 
% team: {\it casc-rose *dot* llnl *dot* gov}.
%    The ROSE project maintain a number of email lists for internal development, 
% external devlopers, and external users. More information is in the ROSE User Manual
% in the Developer's Appendix (see \ref{rose_email_list_info}).

% DQ (1/21/2009): Note this information is maintained in several locations and in LaTeX
% and HTML formats:
%    1) ROSE User Manual: Developer's: Appendix,
%    2) ROSE Installation Manual
%    3) ROSE Tutorial
%    4) API reference web pages for ROSE
%
% DQ (10/28/2008): This information is copied from:
%                  ROSE User Manual: Developer's: Appendix: ROSE Email
%                  Lists.

See https://rosecompiler.org/reference/index.html

%% We have three mailing lists for core developers (those who have write access to
%% the internal repository), all developers (anyone who has write access to the
%% internal or external repository) and all
%% users of ROSE. They are:
%% \begin{itemize}
%% \item rose-core@nersc.gov, web interface: 
%% \htmladdnormallink{https://mailman.nersc.gov/mailman/listinfo/rose-core}{https://mailman.nersc.gov/mailman/listinfo/rose-core}.
%% \item rose-developer@nersc.gov, web interface: 
%% \htmladdnormallink{https://mailman.nersc.gov/mailman/listinfo/rose-developer}{https://mailman.nersc.gov/mailman/listinfo/rose-developer}.
%% \item rose-public@nersc.gov, web interface:
%% \htmladdnormallink{https://mailman.nersc.gov/mailman/listinfo/rose-public}{https://mailman.nersc.gov/mailman/listinfo/rose-public}.
%% \end{itemize}

\subsection{ROSE and the NMI Compile Farm}
\label{NMI_testing}

   The NSF Middleware Initiative (NMI) has provides us with 
time on their system to support the robustness of ROSE across
multiple platforms.  ROSE is not tested on a wide range of platforms
(see table \ref{NMI:prerequisites}).  The prerequisites
used for each platform (machine and operating system) are generated
in the table from the input test descriptions located in the 
directory {\tt ROSE/scripts/nmiBuildAndTestFarm/build\_configs}.

For More information about NMI, see
\htmladdnormallink{{\tt http://nmi.cs.wisc.edu/}}{http://nmi.cs.wisc.edu/}.
To see the details of the ROSE nightly tests click on the link: {\em Run Results}
and select the project, {\em rose compiler}, from the pull down menu.

% This is too long and does not process using LaTeX properly.
% \htmladdnormallink{http://nmi-web.cs.wisc.edu/nmi/}{http://nmi-web.cs.wisc.edu/nmi/}.

   NMI OS and machine (platform) Prerequisites for ROSE:
% Do this when processing latex to generate non-html (not using latex2html)
{\indent
\begin{latexonly}
\begin{figure}[tb]
\begin{center}
\begin{tabular}{|c|} \hline
     {\tt NMI Platform (OS and Machine) Prerequisites}
\\\hline\hline
{
%\footnotesize
\scriptsize
   \lstinputlisting{rosePlatformPrerequisites.txt}
}
\\\hline
\end{tabular}
\end{center}
\caption{ Example NMI machine preques used for nightly tests. }
\end{figure}
\end{latexonly}
\label{NMI:prerequisites}
% End of scope in indentation
}

   NMI OS and machine (platform) Configure Options for ROSE:
% Do this when processing latex to generate non-html (not using latex2html)
{\indent
\begin{latexonly}
\begin{figure}[tb]
\begin{center}
\begin{tabular}{|c|} \hline
     {\tt NMI Platform (OS and Machine) Configure Options}
\\\hline\hline
{
%\footnotesize
\scriptsize
   \lstinputlisting{rosePlatformConfigureOptions.txt}
}
\\\hline
\end{tabular}
\end{center}
\caption{ Example NMI machine configure options used for nightly tests. }
\end{figure}
\end{latexonly}
\label{NMI:configureOptions}
% End of scope in indentation
}

\subsection{Installation Details for Specific Platforms}

See http://rosecompiler.org/ROSE\_HTML\_Reference/installation.html

%%    These are comments collected from users about what 
%% platform specific details they have discovered and shared.
\subsubsection{Fedora 11}
   Ensure required system development packages are installed for your toolchain and
dependencies (e.g., Java runtime for OFP).

\subsubsection{Intel C++ Compiler}
The Intel compiler can run out of space compiling some of the larger files in ROSE.
Although not previously seen by anyone on the ROSE team, one user has reported
that the Intel compiler option {\em -override-limits} was required.  As used on
the following configure line:
More information is on this option and when to use it is at:
{\em http://software.intel.com/en-us/articles/internal-threshold-was-exceeded}.
The problem that this appears to fix is that on some machines the ROSE 
file AST\_FILE\_IO.C will fail with the error {\em memory limit error},
this flag to the Intel compiler will fix this (the mcpcom process goes above 
2.5g memory for this file).


\subsection{Options for Static vs. Dynamic Linking of Executables}
\label{linking_options}

   ROSE supports both static and dynamic linking. Some developers want the
space savings of dynamic linking (also might link faster) and some
want the simplicity of static linked executables (since they can be
easily moved in a single step). Control this with CMake's
{\tt BUILD_SHARED_LIBS} option:
\begin{itemize}
     \item build shared libraries: {\bf -DBUILD\_SHARED\_LIBS=ON}
     \item build static libraries: {\bf -DBUILD\_SHARED\_LIBS=OFF}
\end{itemize}



\subsection{Options to Control the Size of ROSE Executables}

   ROSE by default turns on the compiler's generation of debugging information
(symbol tables and dwarf2 debug information). 
This is done to support development which will nearly always be easier with
the symbol tables generated with the debugging options (typically ``-g'').
However, the volume of code in ROSE will cause a lot of debugging information
to be generated (around 150Meg, just for the debugging information in a 
statically linked executable).

To turn off the generation of debug information in the executables use the
configure options: {\tt --with-CXX\_DEBUG=no --with-C\_DEBUG=no}  Any other
values will use those explicit value to turn on debugging information in
the compilation of ROSE.

The sizes of executables built using dynamic linking are always trivially small,
the significant different matters when executables are built using static linking
(see section~\ref{linking_options}).  For a statically linked executable 
(ROSE-based tool) the size with symbol tables (debugging information) can be
150-200 Meg per executable; this value is for the g++ compiler, other compilers
will vary in the sizes of the executables they generate.  Turning off the 
debug information generated by g++ will reduce the size of the same executable
to be about 40 Meg.  Executables can be stripped of symbol table information
further using the {\tt strip} utility (Linux).  Executables using ROSE that are stripped
will be about 25\% smaller (about 30 Meg).

   Further work in ROSE could likely reduce the size of ROSE based executables,
but it is the judgment of the ROSE team that current work work is sufficient to
generate small enough executables.  If users have need for smaller executables
for ROSE based tools, please contact the ROSE team directly.
